\documentclass[12pt,a4paper]{article}
\usepackage[utf8]{inputenc}
\usepackage{xeCJK}
\usepackage{geometry}
\usepackage{tikz}
\usepackage{flowchart}
\usepackage{listings}
\usepackage{xcolor}
\usepackage{hyperref}
\usepackage{enumitem}

\geometry{margin=2.5cm}
\setCJKmainfont{Noto Sans CJK TC}

% 設定程式碼樣式
\lstset{
    basicstyle=\ttfamily\small,
    backgroundcolor=\color{gray!10},
    frame=single,
    breaklines=true,
    showstringspaces=false
}

\title{\textbf{自動化ICU教學文件生成系統}}
\author{基於Teams錄影轉錄與AI整理}
\date{\today}

\begin{document}

\maketitle

\section{系統概述}
本系統旨在建立一個完整的自動化流程，將ICU住院醫師訓練課程的Teams錄影，透過語音轉錄、AI整理，最終生成結構化的LaTeX教學文件。

\section{流程架構圖}

\begin{center}
\begin{tikzpicture}[node distance=2.5cm, auto]
    % 定義節點樣式
    \tikzstyle{process} = [rectangle, rounded corners, minimum width=3cm, minimum height=1cm, text centered, draw=black, fill=blue!20]
    \tikzstyle{decision} = [diamond, minimum width=3cm, minimum height=1cm, text centered, draw=black, fill=orange!20]
    \tikzstyle{io} = [trapezium, trapezium left angle=70, trapezium right angle=110, minimum width=3cm, minimum height=1cm, text centered, draw=black, fill=green!20]
    \tikzstyle{arrow} = [thick,->,>=stealth]

    % 流程節點
    \node (start) [io] {Teams線上課程};
    \node (record) [process, below of=start] {錄影設定\\繁體中文轉錄};
    \node (transcript) [process, below of=record] {系統自動生成\\中文轉錄文字};
    \node (check) [decision, below of=transcript] {轉錄品質\\檢查};
    \node (manual) [process, right of=check, xshift=3cm] {手動修正\\錯誤轉錄};
    \node (ai) [process, below of=check] {AI文本整理\\結構化處理};
    \node (latex) [process, below of=ai] {生成LaTeX\\教學文件};
    \node (output) [io, below of=latex] {Overleaf編譯\\最終文件};

    % 流程箭頭
    \draw [arrow] (start) -- (record);
    \draw [arrow] (record) -- (transcript);
    \draw [arrow] (transcript) -- (check);
    \draw [arrow] (check) -- node[anchor=west] {需要修正} (manual);
    \draw [arrow] (manual) |- (ai);
    \draw [arrow] (check) -- node[anchor=east] {品質良好} (ai);
    \draw [arrow] (ai) -- (latex);
    \draw [arrow] (latex) -- (output);
\end{tikzpicture}
\end{center}

\section{詳細步驟說明}

\subsection{第一階段：錄影與轉錄}
\begin{enumerate}[itemsep=0.5em]
    \item \textbf{Teams會議設定}
    \begin{itemize}
        \item 開啟會議錄製功能
        \item 語言設定：繁體中文（台灣）
        \item 確保音質清晰，減少背景雜音
    \end{itemize}
    
    \item \textbf{自動轉錄配置}
    \begin{itemize}
        \item 啟用即時轉錄功能
        \item 設定轉錄語言為繁體中文
        \item 錄影結束後自動生成轉錄檔案
    \end{itemize}
\end{enumerate}

\subsection{第二階段：文本預處理}
\begin{enumerate}[itemsep=0.5em]
    \item \textbf{轉錄品質檢查}
    \begin{itemize}
        \item 檢查醫學專業術語正確性
        \item 修正常見轉錄錯誤（如：呼吸器 → 人工呼吸器）
        \item 補充標點符號和段落結構
    \end{itemize}
    
    \item \textbf{內容結構化}
    \begin{itemize}
        \item 識別主要教學重點
        \item 分離理論說明與臨床案例
        \item 標記重要概念和關鍵詞
    \end{itemize}
\end{enumerate}

\subsection{第三階段：AI處理與整理}
\begin{enumerate}[itemsep=0.5em]
    \item \textbf{AI提示詞範例}
    \begin{lstlisting}[language=bash, caption=AI處理指令]
請將以下ICU教學錄音轉錄內容整理成結構化的教學文件：

1. 提取主要教學目標
2. 整理醫學概念和定義
3. 分類臨床案例和處置方式
4. 生成學習重點摘要
5. 輸出為LaTeX格式，包含：
   - 章節標題
   - 醫學術語定義
   - 臨床流程圖
   - 參考文獻格式

內容：[轉錄文字]
    \end{lstlisting}
    
    \item \textbf{AI輸出要求}
    \begin{itemize}
        \item 結構化LaTeX代碼
        \item 醫學術語正確性驗證
        \item 臨床指引符合性檢查
        \item 教學邏輯清晰度評估
    \end{itemize}
\end{enumerate}

\subsection{第四階段：LaTeX文件生成}

\subsubsection{文件模板結構}
\begin{lstlisting}[language=TeX, caption=LaTeX教學文件模板]
\documentclass[12pt,a4paper]{article}
\usepackage[utf8]{inputenc}
\usepackage{xeCJK}
\usepackage{amsmath}
\usepackage{graphicx}
\usepackage{booktabs}
\usepackage{hyperref}

\setCJKmainfont{Noto Sans CJK TC}

\title{ICU住院醫師訓練課程}
\subtitle{[課程主題]}
\author{[講師姓名]}
\date{[錄製日期]}

\begin{document}
\maketitle
\tableofcontents

\section{學習目標}
% AI生成的學習目標

\section{核心概念}
% 醫學概念和定義

\section{臨床應用}
% 實際案例和處置流程

\section{重點摘要}
% 課程重點整理

\section{延伸閱讀}
% 相關參考文獻

\end{document}
\end{lstlisting}

\section{技術實施建議}

\subsection{自動化工具鏈}
\begin{enumerate}
    \item \textbf{Teams API整合}：自動下載轉錄檔案
    \item \textbf{Python腳本}：文本預處理和格式化
    \item \textbf{AI API調用}：ChatGPT/Claude等進行內容整理
    \item \textbf{LaTeX編譯}：Overleaf API或本地編譯環境
    \item \textbf{版本控制}：Git管理文件版本和修訂歷史
\end{enumerate}

\subsection{品質控制機制}
\begin{itemize}
    \item 醫學專業術語數據庫比對
    \item 臨床指引一致性檢查
    \item 教學內容完整性評估
    \item 同儕審核流程整合
\end{itemize}

\section{預期效益}

\begin{itemize}
    \item \textbf{效率提升}：減少90\%手動整理時間
    \item \textbf{標準化}：統一教學文件格式和結構
    \item \textbf{可追溯性}：完整保存教學內容和修訂歷程
    \item \textbf{知識管理}：建立可搜索的教學資源庫
    \item \textbf{持續改進}：透過AI學習優化整理品質
\end{itemize}

\section{結論}

本自動化系統將大幅提升ICU教學文件的生產效率和品質一致性，讓醫師能夠專注於臨床教學本身，而非繁瑣的文件整理工作。透過AI輔助和LaTeX排版，可以產出專業、結構化的教學資源，有效支持住院醫師的學習需求。

\end{document}